\chapter{Polynomials}\label{ch:b1:poly}

Polynomials are the algebraist's favorite functions --- except that
they are not treated as functions here, but as formal expressions in
an indeterminate $X$, added and multiplied by the rules of a
commutative ring. The theory runs strikingly parallel to
\cref{ch:b1:arith}: a Euclidean division, a gcd and Bézout relations,
irreducible elements and a unique factorization. Throughout, $K$
denotes $\Q$, $\R$ or $\C$.

\section{The ring $K[X]$}

\begin{definition}[Polynomial, degree]\label{def:b1:poly:def}
A \emph{polynomial}\index{polynomial} with coefficients in $K$ is a
formal sum
\[
P = a_0 + a_1 X + a_2 X^2 + \dots + a_n X^n
= \sum_{k} a_k X^k,
\]
with $a_k \in K$ all zero from some index on. With the natural
addition and the product
\[
\Bigl(\sum_i a_i X^i\Bigr)\Bigl(\sum_j b_j X^j\Bigr)
= \sum_k \Bigl(\sum_{i+j=k} a_i b_j\Bigr) X^k,
\]
the set $K[X]$ is a commutative ring. The \emph{degree} $\deg P$ of
$P \neq 0$ is the largest $n$ with $a_n \neq 0$; $a_n$ is the
\emph{leading coefficient} ($P$ is \emph{monic}\index{monic
polynomial} when $a_n = 1$), and by convention $\deg 0 = -\infty$.
Every polynomial defines a function $x \mapsto P(x)$ on $K$ by
substitution.
\end{definition}

\begin{proposition}[Degree rules; integral domain]\label{prop:b1:poly:degree}
For $P, Q \in K[X]$:
\[
\deg(P + Q) \leq \max(\deg P, \deg Q),
\qquad
\deg(PQ) = \deg P + \deg Q .
\]
Consequently $K[X]$ is an integral domain, and its units are the
nonzero constants.
\end{proposition}

\begin{proof}
The sum rule is clear (coefficients beyond the max vanish). For the
product, let $a_m$ and $b_n$ be the leading coefficients: the
coefficient of $X^{m+n}$ in $PQ$ is $a_m b_n \neq 0$ ($K$ is a
field, hence an integral domain), and all higher coefficients vanish.
If $P, Q \neq 0$ then $\deg PQ = \deg P + \deg Q \geq 0$, so $PQ \neq
0$: integral domain. If $PQ = 1$ then $\deg P + \deg Q = 0$ forces
$\deg P = \deg Q = 0$: invertible elements are the invertible
constants, i.e.\ all of $K^*$.
\end{proof}

\begin{theorem}[Euclidean division]\label{thm:b1:poly:division}
Let $A, B \in K[X]$ with $B \neq 0$. There is exactly one pair $(Q,
R)$ of polynomials with
\[
A = BQ + R, \qquad \deg R < \deg B .
\]
\end{theorem}

\begin{proof}
\emph{Existence}, by strong induction on $\deg A$. If $\deg A < \deg
B$, take $(Q, R) = (0, A)$. Otherwise write $A = a X^m + \dots$, $B =
b X^n + \dots$ with $m \geq n$; the polynomial $A_1 = A - \frac ab
X^{m-n} B$ has degree $< m$ (the leading terms cancel), so by
induction $A_1 = BQ_1 + R$ with $\deg R < \deg B$, and $A = B(Q_1 +
\frac ab X^{m-n}) + R$.

\emph{Uniqueness}: if $BQ + R = BQ' + R'$, then $B(Q - Q') = R' - R$
with $\deg(R' - R) < \deg B$; by the degree rule this forces $Q - Q'
= 0$, then $R = R'$.
\end{proof}

\begin{example}\label{ex:b1:poly:division}
Divide $A = X^4 + X^3 - 2X + 1$ by $B = X^2 + 1$:
\[
X^4 + X^3 - 2X + 1 = (X^2 + 1)(X^2 + X - 1) + (-3X + 2).
\]
(Compute: subtract $X^2 B$, then $X B$, then $-B$; the remainder
$-3X + 2$ has degree $1 < 2$.)
\end{example}

\begin{remark}[{Arithmetic of $K[X]$}]
With Euclidean division in hand, the entire arithmetic of
\cref{ch:b1:arith} transfers to $K[X]$, with the same proofs, degree
playing the role of absolute value: gcd (normalized to be monic),
extended Euclidean algorithm, Bézout's identity, Gauss's lemma,
irreducible polynomials and unique factorization. We freely use these
transferred results, and \cref{exo:b1:poly:6} rehearses one of them.
\end{remark}

\section{Roots}

\begin{theorem}[Factor theorem]\label{thm:b1:poly:factor}
Let $P \in K[X]$ and $a \in K$. The remainder of $P$ upon division by
$X - a$ is the constant $P(a)$. In particular
\[
P(a) = 0 \iff (X - a) \mid P .
\]
More generally, distinct roots $a_1, \dots, a_r$ of $P$ give the
factorization $P = (X - a_1)\cdots(X - a_r)\, Q$.
\end{theorem}

\begin{proof}
Divide: $P = (X - a) Q + R$ with $\deg R < 1$, so $R$ is a constant
$c$; substituting $X = a$ (substitution respects sums and products)
gives $P(a) = c$. The equivalence follows. For several roots, induct:
$P = (X - a_1) Q_1$, and for $i \geq 2$, $0 = P(a_i) = (a_i - a_1)
Q_1(a_i)$ with $a_i - a_1 \neq 0$ forces $Q_1(a_i) = 0$; apply the
induction hypothesis to $Q_1$.
\end{proof}

\begin{corollary}[A polynomial of degree $n$ has at most $n$ roots]\label{cor:b1:poly:nroots}
A nonzero $P \in K[X]$ of degree $n$ has at most $n$ distinct roots
in $K$. Consequently, a polynomial (of degree $\leq n$) vanishing at
$n + 1$ distinct points is the zero polynomial, and two polynomials
of degree $\leq n$ agreeing at $n+1$ points are equal.
\end{corollary}

\begin{proof}
If $a_1, \dots, a_r$ are distinct roots, \cref{thm:b1:poly:factor}
gives $P = (X-a_1)\cdots(X-a_r) Q$, so $n = \deg P \geq r$. The two
consequences follow by contradiction and by difference.
\end{proof}

\begin{definition}[Derivative, multiplicity]\label{def:b1:poly:derivative}
The \emph{formal derivative} of $P = \sum a_k X^k$ is $P' = \sum_{k
\geq 1} k\,a_k X^{k-1}$; it satisfies the usual rules $(P+Q)' = P' +
Q'$, $(PQ)' = P'Q + PQ'$ (checked on monomials and extended by
linearity). A root $a$ of $P$ has
\emph{multiplicity}\index{multiplicity of a root} $m \geq 1$ when
$(X-a)^m \mid P$ but $(X-a)^{m+1} \nmid P$; the root is \emph{simple}
if $m = 1$, \emph{multiple} if $m \geq 2$.
\end{definition}

\begin{proposition}[Multiplicity via derivatives]\label{prop:b1:poly:multiplicity}
$a$ is a root of $P$ of multiplicity $\geq m$ if and only if
\[
P(a) = P'(a) = \dots = P^{(m-1)}(a) = 0 .
\]
In particular, $a$ is a multiple root of $P$ if and only if $P(a) =
P'(a) = 0$.
\end{proposition}

\begin{proof}
Write $P = (X - a)^m Q + R$ where $R$ is the remainder of the division
by $(X-a)^m$, $\deg R < m$. Differentiating $k \leq m - 1$ times and
evaluating at $a$: the first term contributes $0$ (each derivative
retains a factor $(X-a)$), so $P^{(k)}(a) = R^{(k)}(a)$.

Now a polynomial $R$ of degree $< m$ is determined by $R(a), R'(a),
\dots, R^{(m-1)}(a)$: writing $R = \sum_{k < m} c_k (X - a)^k$
(possible: expand powers of $X = (X - a) + a$), one finds
$R^{(k)}(a) = k!\, c_k$. Hence: all $P^{(k)}(a) = 0$ for $k < m$
$\iff$ all $c_k = 0$ $\iff$ $R = 0$ $\iff$ $(X-a)^m \mid P$.
\end{proof}

\begin{theorem}[Fundamental theorem of algebra]\label{thm:b1:poly:fta}
Every nonconstant polynomial of $\C[X]$ has a root in
$\C$.\index{fundamental theorem of algebra}
\end{theorem}
\admitted

\begin{corollary}[Factorization over $\C$ and over $\R$]\label{cor:b1:poly:factorization}
\begin{enumerate}
  \item Every nonzero $P \in \C[X]$ factors as
        \[
        P = c\, (X - a_1)^{m_1} \cdots (X - a_r)^{m_r},
        \]
        with $c$ the leading coefficient, $a_i$ the distinct complex
        roots, $\sum m_i = \deg P$: \emph{counted with multiplicity,
        a polynomial of degree $n$ has exactly $n$ complex roots}.
  \item Every nonzero $P \in \R[X]$ factors over $\R$ as
        \[
        P = c \prod_i (X - a_i)^{m_i} \prod_j (X^2 + p_j X +
        q_j)^{n_j},
        \]
        the quadratic factors being distinct with $p_j^2 - 4q_j < 0$
        (no real roots).
\end{enumerate}
\end{corollary}

\begin{proof}
(1) Induction on the degree, splitting off one root at a time by
\cref{thm:b1:poly:factor}; the count of degrees matches at each step.

(2) Let $P$ have real coefficients. If $z$ is a complex root of
multiplicity $m$, so is $\conj z$: conjugating $P(z) = 0$ gives
$P(\conj z) = \conj{P(z)} = 0$ (the coefficients are their own
conjugates), and the same applies to the derivatives
(\cref{prop:b1:poly:multiplicity}). Group the non-real roots in
conjugate pairs: each pair contributes
\[
(X - z)(X - \conj z) = X^2 - 2\Re(z)\, X + \abs z^2 ,
\]
a real quadratic with negative discriminant. The real roots contribute
the linear factors.
\end{proof}

\begin{example}\label{ex:b1:poly:factorexample}
$X^4 + 4$ was factored over $\R$ in \cref{exo:b1:complex:5} by pairing
the four complex roots $\pm 1 \pm \iu$:
$X^4 + 4 = (X^2 - 2X + 2)(X^2 + 2X + 2)$. Neither quadratic splits
over $\R$ (discriminants $-4$). Note: an \emph{irreducible} real
polynomial has degree $1$ or $2$ --- this is exactly what the
factorization theorem says.
\end{example}

\section{Coefficients and roots}

\begin{theorem}[Vieta's formulas]\label{thm:b1:poly:vieta}
Let $P = X^n + c_{n-1} X^{n-1} + \dots + c_0$ be monic with roots
$a_1, \dots, a_n \in \C$ (with multiplicity). Then\index{Vieta's
formulas}
\[
\sum_i a_i = -c_{n-1},
\qquad
\sum_{i < j} a_i a_j = c_{n-2},
\qquad \dots, \qquad
a_1 a_2 \cdots a_n = (-1)^n c_0 ,
\]
the $k$-th symmetric function of the roots being $(-1)^k c_{n-k}$.
\end{theorem}

\begin{proof}
Expand $P = (X - a_1)\cdots(X - a_n)$: choosing $X$ in all factors
but $k$ of them produces the term $(-1)^k \bigl(\sum_{i_1 < \dots <
i_k} a_{i_1}\cdots a_{i_k}\bigr) X^{n-k}$; identify coefficients.
\end{proof}

\begin{example}\label{ex:b1:poly:vieta}
For the quadratic $X^2 - sX + p$: sum of roots $s$, product $p$ ---
already used repeatedly (\cref{exo:b1:complex:8}). For a monic cubic
$X^3 + aX^2 + bX + c$ with roots $\alpha, \beta, \gamma$:
\[
\alpha + \beta + \gamma = -a,
\quad
\alpha\beta + \beta\gamma + \gamma\alpha = b,
\quad
\alpha\beta\gamma = -c ,
\]
which allows computing symmetric quantities like $\alpha^2 + \beta^2
+ \gamma^2 = a^2 - 2b$ without solving.
\end{example}

\begin{theorem}[Lagrange interpolation]\label{thm:b1:poly:lagrange}
Let $x_0, \dots, x_n$ be distinct points of $K$ and $y_0, \dots, y_n
\in K$. There is exactly one $P \in K[X]$ of degree $\leq n$ with
$P(x_i) = y_i$ for all $i$, namely\index{Lagrange interpolation}
\[
P = \sum_{i=0}^{n} y_i\, L_i,
\qquad
L_i = \prod_{j \neq i} \frac{X - x_j}{x_i - x_j} .
\]
\end{theorem}

\begin{proof}
Each $L_i$ has degree $n$ and satisfies $L_i(x_i) = 1$, $L_i(x_j) =
0$ for $j \neq i$ (each factor vanishes at the corresponding $x_j$).
So the displayed $P$ has degree $\leq n$ and interpolates. Uniqueness:
two interpolating polynomials of degree $\leq n$ agree at the $n+1$
points $x_i$, hence are equal (\cref{cor:b1:poly:nroots}).
\end{proof}

\section{Exercises}

\begin{exercise}[$\star$]\label{exo:b1:poly:1}
Carry out the Euclidean divisions: $X^5 - 1$ by $X^2 + X + 1$; then
$2X^4 + X^3 - X + 3$ by $X^2 - 2$.
\end{exercise}

\begin{exercise}[$\star$]\label{exo:b1:poly:2}
For which $n \in \N$ does $X^2 + X + 1$ divide $X^{2n} + X^n + 1$?
\emph{Hint: the roots of $X^2 + X + 1$ are $j$ and $j^2$ with $j =
\eu^{2\iu\pi/3}$; discuss $n$ mod $3$.}
\end{exercise}

\begin{exercise}[$\star$]\label{exo:b1:poly:3}
Determine the real $a, b$ so that $(X-1)^2$ divides $P = X^4 + aX^3 +
bX^2 + 1$, then factor $P$ over $\R$ for these values.
\end{exercise}

\begin{exercise}[$\star$]\label{exo:b1:poly:4}
Factor over $\C$ and over $\R$: $X^3 - 1$; $\;X^4 + X^2 + 1$;
$\;X^6 - 1$.
\end{exercise}

\begin{exercise}[$\star\star$]\label{exo:b1:poly:5}
Let $P = X^3 - 6X^2 + 11X - 6$.
\begin{enumerate}
  \item Find the rational roots \emph{(a rational root $p/q$ in
        lowest terms of a monic integer polynomial is an integer
        dividing the constant term --- prove it)}, and factor $P$.
  \item Without solving, compute the sum of the squares and the sum
        of the inverses of the roots via Vieta, and check on the
        factorization.
\end{enumerate}
\end{exercise}

\begin{exercise}[$\star\star$]\label{exo:b1:poly:6}
Compute $\gcd(X^4 - 1,\; X^3 - X^2 + X - 1)$ by the Euclidean
algorithm, and write it as a combination $AU + BV$ of the two
polynomials.
\end{exercise}

\begin{exercise}[$\star\star$]\label{exo:b1:poly:7}
Let $P \in \R[X]$ with $P(x) \geq 0$ for all $x \in \R$. Prove that
$P$ is a sum of two squares of real polynomials: $P = A^2 + B^2$.
\emph{Hint: in the real factorization, real roots have even
multiplicity; write the quadratic factors as $(X - z)(X - \conj z)$
and use $\abs{\,\cdot\,}^2 = (\Re)^2 + (\Im)^2$ on the product of the
$(X - z)$'s.}
\end{exercise}

\begin{exercise}[$\star\star$]\label{exo:b1:poly:8}
Find the polynomial $P$ of degree $\leq 2$ with $P(0) = 1$, $P(1) =
3$, $P(2) = 2$, first by Lagrange's formula, then by solving the
linear system on the coefficients. Verify both answers agree.
\end{exercise}

\begin{exercise}[$\star\star$]\label{exo:b1:poly:9}
Prove that $P = X^{2n+1} - 1$ has exactly one real root, and that for
every $n \geq 1$ the polynomial $1 + X + \frac{X^2}{2!} + \dots +
\frac{X^n}{n!}$ has no multiple root \emph{(compare $P$ and
$P'$)}.
\end{exercise}

\begin{exercise}[$\star\star\star$]\label{exo:b1:poly:10}
(Chebyshev polynomials) Define $T_0 = 1$, $T_1 = X$ and $T_{n+1} =
2X\,T_n - T_{n-1}$.
\begin{enumerate}
  \item Prove by induction that $T_n(\cos\theta) = \cos n\theta$ for
        all $\theta$.
  \item Deduce the $n$ roots of $T_n$ and its leading coefficient.
  \item Prove that $\sup_{x \in \intcc{-1}{1}} \abs{T_n(x)} = 1$,
        attained at $n + 1$ points of $\intcc{-1}{1}$.
\end{enumerate}
\end{exercise}

\begin{exercise}[$\star\star\star$]\label{exo:b1:poly:11}
Let $P \in \C[X]$ be nonconstant with distinct roots $a_1, \dots,
a_r$ (multiplicities $m_1, \dots, m_r$). Prove the identity of
rational functions
\[
\frac{P'(X)}{P(X)} = \sum_{i=1}^{r} \frac{m_i}{X - a_i},
\]
and deduce the Gauss--Lucas theorem: every root of $P'$ lies in the
convex hull of the roots of $P$ \emph{(evaluate the identity at a
root $w$ of $P'$ which is not a root of $P$, take conjugates, and
read the result as $w$ being a weighted average of the $a_i$)}.
\end{exercise}
